%%%%%%%%%%%%%%%%% Оформление ГОСТА%%%%%%%%%%%%%%%%%

% Все параметры указаны в ГОСТЕ на 2021, а именно:

% Шрифт для курсовой Times New Roman, размер – 14 пт.
\setdefaultlanguage[spelling=modern]{russian}
    \setotherlanguage{english}
    
\setmonofont{Times New Roman}
\setmainfont{Times New Roman} 
\setromanfont{Times New Roman} 
\newfontfamily\cyrillicfont{Times New Roman}




% шрифт для URL-ссылок
\urlstyle{same} 

% Междустрочный интервал должен быть равен 1.5 сантиметра.
\linespread{1.5} % междустрочный интервал


% Каждая новая строка должна начинаться с отступа равного 1.25 сантиметра.
\setlength{\parindent}{1.25cm} % отступ для абзаца
\setlist[itemize]{
    leftmargin=0pt,       % Общий отступ всего списка
    labelindent=1.25cm,   % Отступ МЕТКИ от левого края
    itemindent=*,         % Текст начинается под меткой
    labelwidth=1em,       % Ширина области для метки
    labelsep=0.5em,       % Расстояние между меткой и текстом                % Убираем лишние пробелы
    itemsep=0pt,
    topsep=0pt,
    parsep=0pt,
    partopsep=0pt,
    itemsep=0em
    % label=-
}
\setlist[enumerate]{
    leftmargin=0pt,       % Общий отступ всего списка
    labelindent=1.25cm,   % Отступ МЕТКИ от левого края
    itemindent=*,         % Текст начинается под меткой
    labelwidth=1em,       % Ширина области для метки
    labelsep=0.5em,       % Расстояние между меткой и текстом
    itemsep=0pt,          % Убираем лишние пробелы
    topsep=0pt,
    parsep=0pt,
    partopsep=0pt,
    % label=-
}

% Текст, который является основным содержанием, должен быть выровнен по ширине по умолчанию включен из-за типа документа в main.tex


%%%%%%%%%%%%%%%%%% Дополнения %%%%%%%%%%%%%%%%%%%%%%%%%%%%%%%%%

% Путь до папки с изображениями
\graphicspath{ {./Images/} }

% Внесение titlepage в учёт счётчика страниц
\makeatletter
\renewenvironment{titlepage} {
	\thispagestyle{empty}
}


% Цвет гиперссылок и цитирования
\usepackage{hyperref} 
 \hypersetup{ 
     colorlinks=true, 
     linkcolor=black, 
     filecolor=blue, 
     citecolor = black,       
     urlcolor=blue, 
     }
    
\usepackage{caption}

\renewcommand{\figurename}{Рисунок}

\captionsetup[figure]{
  format=plain,
  justification=centering,
  labelsep=none,
  name=Рисунок
}

\makeatletter
\renewcommand{\fnum@figure}{\figurename~\thefigure~–\ } % ← добавлен \ после тире
\makeatother
% Нумерация рисунков
\counterwithout{figure}{section}

% Нумерация таблиц
\counterwithout{table}{section}

\counterwithout{equation}{section}

% шрифт для листингов с лигатурами
\setmonofont{FiraCode-Regular.otf}[
	SizeFeatures={Size=10},
	Path = Settings/,
	Contextuals=Alternate
]

% Перенос текста при переполнении
\emergencystretch=25pt


% настройка подсветки кода и окружения для листингов
%\usemintedstyle{colorful} % делает подсветку для кода
\newenvironment{code}{\captionsetup{type=listing}}{}

% \numberwithin{equation}{section} % если хочешь так же и уравнения нумеровать
\newtheorem{theorem}{Теорема}
\newtheorem{statement}[theorem]{Утверждение}
\newtheorem{lemma}[theorem]{Лемма}
\newtheorem{corollary}{Следствие}
\newtheorem{proposition}[theorem]{Предложение}
\renewcommand{\proofname}{Доказательство}
\newtheorem{example}{Пример}
\newtheorem{definition}{Определение}
% ------------------------------------------
% Настройка выравнивания заголовков
% ------------------------------------------
\usepackage{titlesec}

% Заголовок 1 уровня — по центру, размер шрифта 14pt
\titleformat{\section}
  {\normalfont\bfseries\fontsize{14}{16}\selectfont\centering}
  {\thesection}{1em}{}

% Заголовок 2 уровня — по центру, размер шрифта 14pt
\titleformat{\subsection}
  {\normalfont\bfseries\fontsize{14}{16}\selectfont\centering}
  {\thesubsection}{1em}{}

% Заголовок 3 уровня — по ширине с отступом 1.25 см, размер шрифта 14pt
\titleformat{\subsubsection}
  {\normalfont\bfseries\fontsize{14}{16}\selectfont}
  {\thesubsubsection}{1em}{}
\titlespacing*{\subsubsection}{1.25cm}{3.25ex plus 1ex minus .2ex}{0pt}
\bibpunct{[}{]}{,}{a}{,}{,}

\usepackage{tocloft}

% Добавляем точки ко всем уровням
\renewcommand{\cftsecleader}{\cftdotfill{\cftdotsep}}     % секции
\renewcommand{\cftsubsecleader}{\cftdotfill{\cftdotsep}}  % подсекции
\renewcommand{\cftsubsubsecleader}{\cftdotfill{\cftdotsep}} % подподсекции


% Убираем полужирность у заголовков разделов
\renewcommand{\cftsecfont}{\normalfont}
\renewcommand{\cftsubsecfont}{\normalfont}
\renewcommand{\cftsubsubsecfont}{\normalfont}

% Убираем стандартные вертикальные отступы между разделами
\setlength{\cftbeforesecskip}{0.5em}        % одинаковый отступ перед section
\setlength{\cftbeforesubsecskip}{0.5em}     % одинаковый отступ перед subsection
\setlength{\cftbeforesubsubsecskip}{0.5em}

\renewcommand{\cftsecpagefont}{\normalfont}
\renewcommand{\cftsubsecpagefont}{\normalfont}
\renewcommand{\cftsubsubsecpagefont}{\normalfont}
